\documentclass[../../main.tex]{subfiles}% 注意这里的文件路径不能用 ./main.tex ,否则用latexmk编译子文件会报错
\graphicspath{{\subfix{./image/}}} % 指定图片目录,后续可以直接使用图片文件名
% 注意这里的文件路径不能用 ../../image/ ,否则用latexmk编译子文件会报错

% 例如:
% \begin{figure}[H]
% \centering
% \includegraphics[scale=0.3]{图.png}
% \caption{}
% \label{figure:图}
% \end{figure}
% 注意:上述\label{}一定要放在\caption{}之后,否则引用图片序号会只会显示??.

\begin{document}

\section{曲线的弧长}

正则曲线的弧长是刚体运动不变量，与参数选取无关。
\begin{definition}[曲线的弧长]
设正则参数曲线$\bm{r}(t)$，$t\in[a,b]$，其**弧长**定义为
\[
s = \int_a^b \big|\bm{r}'(t)\big|\mathrm{d}t.
\]
其中$\big|\bm{r}'(t)\big|=\sqrt{x'(t)^2+y'(t)^2+z'(t)^2}$为切向量的长度。
\end{definition}

弧长的微分形式称为**弧长元素**：
\[
\mathrm{d}s = \big|\bm{r}'(t)\big|\mathrm{d}t,\quad \left(\dfrac{\mathrm{d}s}{\mathrm{d}t}\right)^2 = \bm{r}'(t)\cdot\bm{r}'(t).
\]


\begin{theorem}
正则曲线$\bm{r}(t)$的参数$t$为**弧长参数**（记为$s$）的充要条件是
\[
\big|\bm{r}'(s)\big| \equiv 1.
\]
\end{theorem}

弧长参数的意义：切向量为单位向量，参数本身具有几何不变性，是曲线的自然参数。

\begin{remark}
弧长参数仅差一个常数（起点选取），是曲线最自然的参数形式，后续理论均以弧长为参数展开。
\end{remark}

\begin{exercise}
\begin{enumerate}
\item 求曲线$\bm{r}(t)=(\cosh t, \sinh t, t)$在$[-1,1]$上的弧长。
\item 求曳物线$\bm{r}(t)=(\cos t, \ln(\sec t+\tan t)-\sin t)$从$0$到$t_0$的弧长。
\end{enumerate}
\end{exercise}






\end{document}